% This file should be replaced with your file with thesis content.
%=========================================================================
% Authors: Michal Bidlo, Bohuslav Křena, Jaroslav Dytrych, Petr Veigend and Adam Herout 2019

% For compilation piecewise (see projekt.tex), it is necessary to uncomment it and change
% \documentclass[../projekt.tex]{subfiles}
% \begin{document}

\chapter{Introduction}

@TODO: after finishing the thesis

úvod, asi trochu k design philosophy a cílům

The goal of this thesis is the design and implementation of a dynamic injury system of game characters in the Godot Engine. The term \emph{character} here describes game characters in general, and is not intended to be limited to anthropomorphs.

This demo is not player-centric, but rather aims to implement a standalone injury system and its effects on the creature's behavior and animation. That is not to say player experience and game design is thrown out of the window, only that the player controller is not the focus of this demo.

The Godot Engine was chosen for this project because 

Chapter 2 

% ===================================================================
% ===================================================================
% ===================================================================

\chapter{Game Character Injury Systems and Animation (Theory)}

Combat and injury have been a part of human history even since prehistoric times. Singular fights, epic battles, and strifeful wars riddle our existence in stories -- mythological through modern, in our personal lives, or in large-scale historical real life events. It is then no wonder that games have used this theme in a myriad of different ways. From mythical, heroic combat with swords and magic to competetive shooters, countless games had to contend with designing and implementing combat and injury systems.

Each game has its budget and own set of priorities based on their design. Sometimes, developers may decide to prioritize other parts of the game like visuals, story, and voice acting, over the implementation of complex combat systems, opting for greater abstraction. The easiest systems are the ones that model characters as a singular entity with no subparts and a pool of hit points that determine its binary state -- alive or dead. However, other designs call for a more complex and varied experience. We will look into those in a later section.

In this chapter, we will delve into the basics of game design and the questions that a designer should answer before creating a system for a game. Next, we will look into the subsystems constituting an advanced injury system of game characters -- systems of locomotion (movement), tracking and simulating harm and injury through subsystems, and having enemy AI reflect the injuries in its behavior.

\section{Game Design}
Before designing a game system, we need to familiarize ourselves with the theory, practice, and questions of game design.

Firstly, what is game design? \emph{Game design is the act of deciding what a game should be. Decisions about rules, look and feel, timing, pacing, risk-taking, rewards, punishments, and everything else the player experiences is the responsibility of the game designer} {[cite Schell book of lenses]}.

\emph{Rules are an essential feature of games. A game's outcome should, to a certain extent, be unpredictable. Games that are fully predictable (a predetermined sequence of events) or fully unpredictable (pure games of chance) aren't engaging. Players want their skills and their strategic decisions to make more of a difference, or else they become frustrated. Skill should also contribute to victory. The longer the game is, the more true this is.}

\emph{In addition to chance, choices made by players and complex gameplay created by the rules are other ways to make games unpredictable. Complex systems usually have many interacting parts, often simple rules that can combine into interwoven and. difficult to foresee situations. In this type of game, the ability to read a game's current state and understand its strategic complexities is the most important game-playing skill.}

\emph{Most games mix these three sources of unpredictability. They include an element of chance, player choices, and complex rules.} {[cite Advanced Game Mechanics]}

Now, we can finally get into the core of this thesis and answer the necessary questions about our game design. The goal is to create a more complex combat experience with less abstractions, where the player can make meaningful choices about where to attack an enemy to cripple its capabilities based on those attacks. For example, a player that prefers melee combat might decide to target an enemy's wings first to make it lose its ability to fly to hinder its maneuverability and disable it from getting out of melee range.

Let's dissect the game design from the perspective of the three sources of unpredictability that should make the game more engaging, as mentioned above. The element of chance comes from the enemy's AI which makes randomized decisions based on weighted transitions in its behavior tree. The player chooses which body parts to attack so that they can disable the enemy's abilities that require those parts to be triggered. The complexity comes from the way body parts serve as requisites for the enemy's capabilities, such as two wings being necessary for flight, at least three limbs for walking, or a stinger tail for poison attacks, and the way they enable and limit what actions it has available in its behavior.

We could also argue that depending on the implementation of the player controller, there are even more elements, such as the player maneuvering around the virtual space to avoid or close in on an enemy, or any of the player's special abilities and attacks with various ranges and cooldowns. However, the player controller is a complex issue of its own. It is not covered by this thesis, but the player's experience is a crucial aspect of games and should be accounted for when discussing of game design ([source book of lenses probably?]).

While assessing, designing, and implementing these features, the goal is also to develop in-engine tools that would make creating such enemies more streamlined in Godot.


\section{Systems of Locomotion}
simulating movement, procedural movement and animation, possible techniques (like IK), adjusting based on centroid (e.g., different foot placement and leaning if a leg is gone)

The injury system has to be dynamic, adapting to whichever injuries are applied to the character's body. Hence procedural animation fits this goal of adaptibility well. Rather than traditional animation with preset model movements, procedural animation can account for and adapt to injuries on the run, and make for a more believable and immersive interaction with the injured character.

\subsection{Inverse Kinematics}
\textbf{![TODO]} Describe IK from survey\_IK!

Inverse kinematics is "rotational and translational motion of
points, bodies (objects) and systems of bodies (groups of objects)
without consideration of what causes the motion or any reference
to mass, force or torque"

It had originally been used in robotics for moving arms to targets, but has also expanded into animation in computer graphics. In the latter field, it is very useful for animating \emph{articulated figures} (figures with a specific, solid form like humans, animals, and generally creatures with limbs, as opposed to blobs of jelly, for example). These models are often complicated, and require efficient calculation of complex and multi-link models that can have specified constraints (e.g., a maximum angle that a link can turn by).

[description of IK basic mechanisms and illustrative images - end effectors, end points and target points, joints and links] The multi-link models in inverse kinematics are usually calculated by points (target positions to reach and the model's joints, known as \emph{end effectors}) and their potential constraints. These end effectors are specified by the model (the joints, the links between them, and the constraints limiting their movement) or by the animation system (the target points that a multi-link model, such as an arm, needs to reach to accomplish its task).

A model is a series of links connected by joints. Joints are the main component allowing movement -- their positions determine the positions and orientations of the links. Links are the "connective tissue" between joints and are defined by their length, shape, volume, and mass properties. A \emph{posture} is the skeletal configuration of the figure and must most often satisfy constraints to adhere to realistic movement, such as the torque limitations of each joint or the links not intersecting with other models in their path. A skeletal configuration is usually separated into \emph{chains}, standalone parts attached to a static point on the model (e.g., the limbs of animals and anthropomorphs). The chains usually work in a hierarchy, and the movement of one joint affects the movement of the ones before it.

[illustration image of chained movement effects]

\subsubsection{Usage of Inverse Kinematics}

Forward kinematics (traditional animation where known, predefined transformations are applied to the model) can fall short in situations where we might need to dynamically adjust the target locations, such as characters standing on slopes or an appandage reaching for a target that can be a variable distance or height away. Inverse kinematics is adaptive, the movement of the links (such as limbs) is dynamically calculated and placed, making for natural-looking movement. This quality is what makes it an ideal animation method for our dynamic limb-based injury system.

[image example(s) of IK use in video games where an end effector target - legs touching the actual ground, nathan drake touching walls, dynamic walking systems]


\section{Simulating Harm and Injury}
possible systems - HP-based; body-part based; idk

probably going with a joint HP pool where each body part is some fraction of that pool

the problem with critical injuries and instant death -- should we aim to simulate injury realistically and allow instant death in cases of striking a critical body part like the head, or should we enforce a design decision and disallow instant death to preserve game balance?

\section{Enemy AI}
deciding on actions and movement; limiting options based on harm; procedural animation of actions

\section{Case Studies}
Nastudujte existující a možná řešení dané problematiky -- jak to resi ostatni hry?

Horizon Zero Dawn/Forbidden West
Monster Hunter
Dead Space
Dragon's Dogma
The Surge
Underrail
Crackdown
Vagrant Story
Fallout
Shadow Tower: Abyss
Hybrid Heaven
Wasteland
Bushido Blade
Binary Domain
Phoenix Point
Xenoblade Chronicles X
Mechwarrior
Assassin's Creed Valhalla

Kingdom Come: Deliverance
Fear \& Hunger
Escape from Tarkov
Project Zomboid
Neo Scavenger

% ===================================================================
% ===================================================================
% ===================================================================

\chapter{Injury System Design}
Navrhněte uživatelské rozhraní a výslednou funkcionalitu aplikace, která demonstruje pokročilé reakce herní postavy na zranění.

Design justified by design principles (books).

Godot plugin design and mockups


% ===================================================================
% ===================================================================
% ===================================================================

\chapter{Implementation and Demonstration in Various Test Cases/Situations}
\section{Creating an Enemy Model}
Blender. Explain enemy design (body parts) -- that's why it's a custom model.
sculpting, retopo, rigify addon skeleton edited to fit the model, automatic weight paints with some manual adjustments, exported to Godot.

\section{Inverse Kinematics in Godot}
No need to reinvent the wheel, just make sure it works well. Using Godot 4's SkeletonIK3D and LookAtModifier3D.


% ===================================================================
% ===================================================================
% ===================================================================

\chapter{User/Performance testing of the resulting implemented design}


% ===================================================================
% ===================================================================
% ===================================================================

\chapter{Documentation and Publication}
+ video
